% ============================================================================
%  ShowerShapesAnalysis — Software Documentation
% ============================================================================
\documentclass[12pt,a4paper]{article}

% --- Packages ---------------------------------------------------------------
\usepackage[utf8]{inputenc}
\usepackage[T1]{fontenc}
\usepackage{lmodern}
\usepackage[margin=2.5cm,headheight=15pt]{geometry}
\usepackage{amsmath,amssymb}
\usepackage{graphicx}
\usepackage{booktabs}
\usepackage{caption}
\usepackage{hyperref}
\usepackage{cleveref}
\usepackage{xcolor}
\usepackage{float}
\usepackage{siunitx}
\usepackage{fancyhdr}
\usepackage{listings}
\usepackage{enumitem}
\usepackage{longtable}
\usepackage{array}

% --- Page style -------------------------------------------------------------
\pagestyle{fancy}
\fancyhf{}
\fancyhead[L]{\small ShowerShapesAnalysis — Software Documentation}
\fancyhead[R]{\small ATLAS Internal}
\fancyfoot[C]{\thepage}
\renewcommand{\headrulewidth}{0.4pt}

% --- Hyperref setup ---------------------------------------------------------
\hypersetup{
  colorlinks=true,
  linkcolor=blue!60!black,
  citecolor=blue!60!black,
  urlcolor=blue!60!black
}

% --- Listings style ---------------------------------------------------------
\definecolor{codebg}{gray}{0.96}
\definecolor{codegreen}{rgb}{0.2,0.5,0.2}
\definecolor{codegray}{rgb}{0.5,0.5,0.5}

\lstdefinestyle{shellstyle}{
  backgroundcolor=\color{codebg},
  basicstyle=\ttfamily\small,
  breaklines=true,
  frame=single,
  rulecolor=\color{codegray},
  xleftmargin=4pt,
  xrightmargin=4pt,
  framesep=3pt,
  numbers=none,
  showstringspaces=false,
  commentstyle=\color{codegreen}\itshape,
  keywordstyle=\color{blue!70!black}\bfseries,
  stringstyle=\color{red!60!black},
  language=bash,
  morekeywords={root,python,bash,setupATLAS,lsetup,asetup,pathena,
    voms-proxy-init,source,cd,mkdir},
  literate={~}{{\raisebox{0.5ex}{\texttildelow}}}{1},
}

\lstdefinestyle{cppstyle}{
  backgroundcolor=\color{codebg},
  basicstyle=\ttfamily\small,
  breaklines=true,
  frame=single,
  rulecolor=\color{codegray},
  xleftmargin=4pt,
  xrightmargin=4pt,
  framesep=3pt,
  numbers=left,
  numberstyle=\tiny\color{codegray},
  showstringspaces=false,
  commentstyle=\color{codegreen}\itshape,
  keywordstyle=\color{blue!70!black}\bfseries,
  stringstyle=\color{red!60!black},
  language=C++,
}

\lstset{style=shellstyle}

% --- Title ------------------------------------------------------------------
\title{%
  \textbf{ShowerShapesAnalysis}\\[6pt]
  \large Software Documentation\\[4pt]
  \large Photon Shower Shape Reweighting — MC23e\\[4pt]
  \normalsize Version 1.0
}
\author{ATLAS Simulation}
\date{\today}

% ============================================================================
\begin{document}
\maketitle
\thispagestyle{fancy}

\tableofcontents
\clearpage

% ============================================================================
\section{Overview}
\label{sec:overview}
% ============================================================================

\texttt{ShowerShapesAnalysis} is a ROOT/Python analysis package for
studying and correcting photon electromagnetic shower shape variables
using cell-level calorimeter information from ATLAS MC23e simulated
ntuples.

The package provides two main capabilities:

\begin{enumerate}[leftmargin=2em]
  \item \textbf{Shower shape validation}: Recomputes the shower shape
    variables $R_\eta$, $R_\phi$, and $w_{\eta\,2}$ directly from
    Layer-2 cell energies and compares them against the stored (fudged
    and unfudged) values in the ntuples. This validates the cell-level
    computation.

  \item \textbf{Cell-energy reweighting closure test}: A complete
    pipeline that generates pseudo-data by distorting MC cell energies,
    derives per-cell reweighting corrections, applies them, and
    evaluates how well the corrected MC recovers the pseudo-data
    distributions. Two correction methods are implemented and compared.
\end{enumerate}

The package is designed to run on CERN lxplus (EL9) with ROOT~6.32+
available via the ATLAS software distribution on CVMFS. It operates on
ntuples produced by the \texttt{NTupleMaker} Athena algorithm, which
stores the $7\times11$ Layer-2 cell energies for each photon candidate.

% ============================================================================
\section{Software Requirements}
\label{sec:requirements}
% ============================================================================

\subsection{Operating system}

The package is developed and tested on CERN lxplus nodes running
AlmaLinux~9 (EL9). It should work on any Linux system with access to
CVMFS.

\subsection{Required software}

\begin{table}[htbp]
\centering
\caption{Software dependencies.}
\label{tab:requirements}
\begin{tabular}{llp{7.5cm}}
\toprule
\textbf{Software} & \textbf{Version} & \textbf{Purpose} \\
\midrule
ROOT & $\geq$ 6.32 & Event processing, histogramming, plotting, C++ JIT
  compilation via \texttt{gInterpreter} \\
Python & 3.9+ & Shower shape validation scripts
  (bundled with ROOT distribution) \\
Athena & 25.0.40 & NTupleMaker compilation (ntuple production only) \\
pathena & (via \texttt{lsetup panda}) & Grid job submission
  (ntuple production only) \\
rucio & (via \texttt{lsetup rucio}) & Dataset management
  (ntuple production only) \\
\bottomrule
\end{tabular}
\end{table}

\paragraph{Note.} Athena, pathena, and rucio are required only for
producing the input ntuples from AOD files on the grid. If pre-produced
ntuples are already available, only ROOT~6.32+ is needed.

\subsection{Python libraries}

The Python scripts use only \texttt{ROOT} (PyROOT) and \texttt{numpy},
both of which are included in the ROOT distribution available on CVMFS.
No additional \texttt{pip install} is required.

% ============================================================================
\section{Installation}
\label{sec:installation}
% ============================================================================

\subsection{Obtaining the repository}

Clone the repository from GitHub:

\begin{lstlisting}
git clone git@github.com:martafsilva31/ShowerShapesAnalysis.git
cd ShowerShapesAnalysis
\end{lstlisting}

\subsection{Environment setup}

On lxplus, source the provided setup script. This configures the ATLAS
software environment (Athena + ROOT) and optionally sources the
NTupleMaker build:

\begin{lstlisting}
# Navigate to the parent ShowerShapes directory
cd /path/to/ShowerShapes/
source ShowerShapesAnalysis/setup.sh
\end{lstlisting}

The setup script performs the following steps:
\begin{enumerate}[leftmargin=2em]
  \item Sets environment variables \texttt{SHOWERSHAPES\_DIR} and
    \texttt{ANALYSIS\_DIR}.
  \item Runs \texttt{setupATLAS} and \texttt{asetup Athena,25.0.40}.
  \item Sources the NTupleMaker build (if present).
  \item Defines convenience aliases (\texttt{cdanalysis},
    \texttt{cdntuples}, etc.).
\end{enumerate}

\paragraph{Minimal setup (analysis only).} If ntuples are already
available and no Athena build is needed, a lightweight alternative is:

\begin{lstlisting}
setupATLAS
lsetup "root 6.32.08-x86_64-el9-gcc13-opt"
\end{lstlisting}

\subsection{Compilation}

No compilation step is required. The C++ ROOT macros are compiled
at runtime via ROOT's JIT compiler (\texttt{gInterpreter}). The closure
test macros include \texttt{config.h} as a shared header, which ROOT
processes automatically when the macro is invoked.

% ============================================================================
\section{Repository Structure}
\label{sec:structure}
% ============================================================================

\begin{table}[htbp]
\centering
\caption{Top-level directory layout.}
\label{tab:dirs}
\small
\begin{tabular}{lp{10cm}}
\toprule
\textbf{Path} & \textbf{Description} \\
\midrule
\texttt{scripts/}         & Analysis scripts (Python + ROOT C++ macros) \\
\texttt{scripts/closure\_test/} & Closure test pipeline (6 C++ macros +
  shell drivers) \\
\texttt{grid/}             & Grid submission scripts for ntuple production
  via pathena \\
\texttt{grid/samples/}     & Sample lists (dataset container names) \\
\texttt{configs/}          & Configuration files (reserved for future use) \\
\texttt{ntuples/}          & Input ntuples (gitignored; see \cref{sec:input}) \\
\texttt{output/}           & All generated output: histograms, plots, closure
  test results (gitignored) \\
\texttt{report/}           & \LaTeX{} reports and compiled PDFs \\
\texttt{setup.sh}          & Environment setup script \\
\texttt{.github/}          & Copilot instruction files for AI-assisted
  development \\
\bottomrule
\end{tabular}
\end{table}

\subsection{Validation scripts (\texttt{scripts/})}

\begin{table}[htbp]
\centering
\caption{Shower shape validation scripts.}
\label{tab:scripts:validation}
\small
\begin{tabular}{lp{9cm}}
\toprule
\textbf{File} & \textbf{Description} \\
\midrule
\texttt{utils.py}
  & Shared Python utilities: argument parsing, TTree I/O helpers,
    cell-energy window extraction, $w_{\eta\,2}$ computation. \\
\texttt{compute\_reta.py}
  & Recomputes $R_\eta = E(3\times7)/E(7\times7)$ from cell energies.
    Compares with stored unfudged and fudged values. Outputs histogram
    ROOT file. Uses JIT-compiled C++ for the inner loop. \\
\texttt{compute\_rphi.py}
  & Same as above for $R_\phi = E(3\times3)/E(3\times7)$. \\
\texttt{compute\_weta\_2.py}
  & Same for $w_{\eta\,2}$. Uses cluster $\eta_2$ plus relative cell
    offsets. \\
\texttt{plot\_reta.C}
  & ROOT macro: reads histogram file, produces a single-page PDF with
    an overlay of the three distributions (unfudged, fudged, computed)
    and a ratio panel. \\
\texttt{plot\_rphi.C}
  & Same layout for $R_\phi$. \\
\texttt{plot\_weta\_2.C}
  & Same layout for $w_{\eta\,2}$. \\
\bottomrule
\end{tabular}
\end{table}

\subsection{Closure test pipeline (\texttt{scripts/closure\_test/})}

\begin{table}[htbp]
\centering
\caption{Closure test source files and their roles.}
\label{tab:scripts:closure}
\small
\begin{tabular}{lp{9cm}}
\toprule
\textbf{File} & \textbf{Description} \\
\midrule
\texttt{config.h}
  & Shared header defining cluster geometry, $|\eta|$ binning (14~bins),
    energy-fraction binning (12~bins), branch names, quality-cut
    functions, shower shape formulas, and the distortion bias pattern. \\
\texttt{create\_pseudodata.C}
  & Generates a pseudo-data ROOT file by applying systematic per-cell
    bias plus Gaussian noise to MC cell energies. Conserves total
    cluster energy. \\
\texttt{derive\_corrections.C}
  & Derives reweighting corrections by comparing MC and pseudo-data cell
    fractions. Produces Method~1 (flat shift) and Method~2 (TProfile)
    corrections, stored in a single ROOT file. \\
\texttt{apply\_corrections.C}
  & Reads corrections and applies them to MC cell energies. Outputs a
    reweighted ntuple. Accepts a method flag (1 or 2). \\
\texttt{validate\_closure.C}
  & Fills shower shape histograms ($R_\eta$, $R_\phi$, $w_{\eta\,2}$)
    per $|\eta|$ bin for MC, pseudo-data, Method~1, and Method~2.
    Computes $\chi^2/n_\text{df}$ closure metrics. \\
\texttt{plot\_closure.C}
  & Generates multi-page comparison PDFs (one page per $|\eta|$ bin)
    with overlay + ratio panels. \\
\texttt{run\_closure\_test.sh}
  & Driver script: runs all 6 pipeline steps for a single distortion
    level. \\
\texttt{run\_closure\_suite.sh}
  & Wrapper: runs \texttt{run\_closure\_test.sh} at 1\%, 3\%, and 5\%
    distortion levels. \\
\bottomrule
\end{tabular}
\end{table}

\subsection{Grid submission (\texttt{grid/})}

\begin{table}[htbp]
\centering
\caption{Grid submission scripts.}
\label{tab:scripts:grid}
\small
\begin{tabular}{lp{9cm}}
\toprule
\textbf{File} & \textbf{Description} \\
\midrule
\texttt{submit\_mc23e\_Zllg.sh}
  & Submits pathena jobs for MC23e $Z\to\ell\ell\gamma$ samples
    (DSIDs 700770, 700771). \\
\texttt{submit\_data22.sh}
  & Submits jobs for 2022 data runs. \\
\texttt{submit\_data24.sh}
  & Submits jobs for 2024 data (single run). \\
\texttt{submit\_data24\_all.sh}
  & Submits jobs for all 2024 data runs. \\
\bottomrule
\end{tabular}
\end{table}

% ============================================================================
\section{Input Data}
\label{sec:input}
% ============================================================================

\subsection{Expected directory layout}

The analysis expects input ntuples under \texttt{ntuples/} with the
following structure:

\begin{lstlisting}
ntuples/
  mc23e/
    mc23e_700770_Zeeg.root       # Z -> eeg MC (DSID 700770)
    mc23e_700771_Zmumug.root     # Z -> mumug MC (DSID 700771)
  data22/
    data22_XXXXXXXX_Zeeg.root    # Data 2022 runs
\end{lstlisting}

\subsection{Ntuple format}

The ntuples are produced by the \texttt{NTupleMaker} Athena algorithm
and stored in a TTree named \texttt{tree}. The analysis uses the
following branches:

\begin{table}[htbp]
\centering
\caption{Required TTree branches (\texttt{tree}).}
\label{tab:branches}
\small
\begin{tabular}{lll}
\toprule
\textbf{Branch name} & \textbf{Type} & \textbf{Description} \\
\midrule
\texttt{photon.7x11ClusterLr2E}
  & \texttt{vector<double>}
  & 77 cell energies (MeV) in the $7\eta\times11\phi$ Layer-2 window.
    Stored in row-major order: $\phi$ varies fastest. \\[3pt]
\texttt{photon.7x11ClusterLr2Size}
  & \texttt{int}
  & Number of cells (always 77 for valid clusters). \\[3pt]
\texttt{photon\_cluster.eta2}
  & \texttt{float}
  & Cluster barycentre pseudorapidity in Layer~2. \\[3pt]
\texttt{photon.reta}
  & \texttt{float}
  & Stored $R_\eta$ (fudged value). \\[3pt]
\texttt{photon.unfudged\_reta}
  & \texttt{float}
  & Stored $R_\eta$ (unfudged value). \\[3pt]
\texttt{photon.rphi}
  & \texttt{float}
  & Stored $R_\phi$ (fudged value). \\[3pt]
\texttt{photon.unfudged\_rphi}
  & \texttt{float}
  & Stored $R_\phi$ (unfudged value). \\[3pt]
\texttt{photon.weta2}
  & \texttt{float}
  & Stored $w_{\eta\,2}$ (fudged value). \\[3pt]
\texttt{photon.unfudged\_weta2}
  & \texttt{float}
  & Stored $w_{\eta\,2}$ (unfudged value). \\
\bottomrule
\end{tabular}
\end{table}

\paragraph{Cell energy layout.}
The 77 cell energies are arranged in a $7\times11$
($\eta\times\phi$) grid. The flat index is calculated as:
\begin{equation}
  \text{index} = \eta_\text{row} \times 11 + \phi_\text{col},
  \qquad
  \eta_\text{row} \in [0,6],\;
  \phi_\text{col} \in [0,10].
\end{equation}
The central cell corresponds to $(\eta_\text{row}, \phi_\text{col}) =
(3, 5)$, i.e.\ flat index~38.

\subsection{Sample details}

The primary MC sample used is:
\begin{itemize}[leftmargin=2em]
  \item \textbf{DSID 700770}:
    \texttt{Sh\_2214\_eegamma} ($Z\to ee\gamma$), MC23e 13.6~TeV
  \item $\sim$9.6~million events, $\sim$13~GB
  \item AOD container: \texttt{mc23\_13p6TeV.700770.Sh\_2214\_eegamma.%
    merge.AOD.e8514\_e8586\_s4369\_s4370\_r16083\_r15970}
\end{itemize}

A secondary sample (DSID 700771, $Z\to\mu\mu\gamma$) is available for
cross-checks.

% ============================================================================
\section{Running the Analysis}
\label{sec:running}
% ============================================================================

All commands below assume the environment has been set up as described
in \cref{sec:installation}.

% ---------- 6.1 Validation -------------------------------------------------
\subsection{Shower shape validation}
\label{sec:running:validation}

The validation workflow computes shower shape variables from cell
energies and compares them against stored values. It consists of two
steps per variable: a Python script that fills histograms, and a ROOT
macro that produces comparison plots.

\subsubsection{Step 1: Compute histograms}

Run the Python compute scripts from the \texttt{scripts/} directory.
Each script takes an input ntuple and produces a ROOT file containing
histograms:

\begin{lstlisting}
cd scripts/

# R_eta
python compute_reta.py \
  -i ../ntuples/mc23e/mc23e_700770_Zeeg.root \
  -o ../output/reta_Zeeg.root

# R_phi
python compute_rphi.py \
  -i ../ntuples/mc23e/mc23e_700770_Zeeg.root \
  -o ../output/rphi_Zeeg.root

# w_eta_2
python compute_weta_2.py \
  -i ../ntuples/mc23e/mc23e_700770_Zeeg.root \
  -o ../output/weta2_Zeeg.root
\end{lstlisting}

Repeat with the $Z\to\mu\mu\gamma$ ntuple (replace
\texttt{mc23e\_700770\_Zeeg.root} with
\texttt{mc23e\_700771\_Zmumug.root} and adjust output names
accordingly).

\paragraph{Parameters.}
\begin{itemize}[leftmargin=2em]
  \item \texttt{-i}, \texttt{-{}-input}: Path to input ROOT ntuple
    (required).
  \item \texttt{-o}, \texttt{-{}-output}: Path to output ROOT
    histogram file (required).
  \item \texttt{-n}, \texttt{-{}-max-events}: Maximum number of events
    to process. Default: $-1$ (all events).
\end{itemize}

\subsubsection{Step 2: Generate plots}

Run the ROOT plotting macros. Each macro reads the histogram file from
Step~1 and produces a single-page PDF:

\begin{lstlisting}
# R_eta plot
root -l -b -q \
  'plot_reta.C("../output/reta_Zeeg.root","../output/plots/reta_Zeeg")'

# R_phi plot
root -l -b -q \
  'plot_rphi.C("../output/rphi_Zeeg.root","../output/plots/rphi_Zeeg")'

# w_eta_2 plot
root -l -b -q \
  'plot_weta_2.C("../output/weta2_Zeeg.root","../output/plots/weta2_Zeeg")'
\end{lstlisting}

The output argument is the base filename (without extension); the macro
appends \texttt{.pdf} automatically. Create the output directory first
if it does not exist:

\begin{lstlisting}
mkdir -p ../output/plots
\end{lstlisting}

Each plot contains:
\begin{itemize}[leftmargin=2em]
  \item \textbf{Upper panel}: Normalised distributions of three
    variants — Unfudged/stored (blue), Fudged/stored (red), Computed
    from cells (green dashed).
  \item \textbf{Lower panel}: Ratio to the Unfudged baseline
    (range 0.8--1.2).
\end{itemize}

% ---------- 6.2 Closure test -----------------------------------------------
\subsection{Closure test}
\label{sec:running:closure}

The closure test validates the cell-energy reweighting method using
MC-only data. The pipeline consists of six sequential stages.

\subsubsection{Automated execution}

\paragraph{Single distortion level.}

\begin{lstlisting}
cd scripts/closure_test/

# Run with 500k events at 5% distortion
bash run_closure_test.sh 500000 0.05
\end{lstlisting}

\paragraph{Full suite (1\%, 3\%, 5\%).}

\begin{lstlisting}
# Run all three distortion levels
bash run_closure_suite.sh 500000
\end{lstlisting}

For full-statistics runs using all $\sim$9.6M events:

\begin{lstlisting}
bash run_closure_suite.sh -1
\end{lstlisting}

\subsubsection{Command-line parameters}

\paragraph{\texttt{run\_closure\_test.sh}} accepts three positional
arguments:

\begin{table}[htbp]
\centering
\caption{Parameters for \texttt{run\_closure\_test.sh}.}
\label{tab:params:closure}
\begin{tabular}{llll}
\toprule
\textbf{Position} & \textbf{Name} & \textbf{Default} &
  \textbf{Description} \\
\midrule
1 & \texttt{N\_EVENTS} & 500000 & Number of events ($-1$ = all) \\
2 & \texttt{DIST\_LEVEL} & 0.05 & Systematic distortion magnitude \\
3 & \texttt{NOISE\_SIGMA} & $-1$ & Per-cell noise sigma
  ($-1$ = auto: $\delta/2$) \\
\bottomrule
\end{tabular}
\end{table}

\paragraph{\texttt{run\_closure\_suite.sh}} accepts one positional
argument:

\begin{table}[htbp]
\centering
\caption{Parameters for \texttt{run\_closure\_suite.sh}.}
\label{tab:params:suite}
\begin{tabular}{llll}
\toprule
\textbf{Position} & \textbf{Name} & \textbf{Default} &
  \textbf{Description} \\
\midrule
1 & \texttt{N\_EVENTS} & 500000 & Events per level ($-1$ = all) \\
\bottomrule
\end{tabular}
\end{table}

The suite always runs at distortion levels $\delta = 0.01, 0.03, 0.05$.

\subsubsection{Manual step-by-step execution}

Each step can be executed individually for debugging or partial reruns.
All commands assume the working directory is
\texttt{scripts/closure\_test/} and use the following variables:

\begin{lstlisting}
INPUT=../../ntuples/mc23e/mc23e_700770_Zeeg.root
OUTDIR=../../output/closure_test_5pct
LEVEL=0.05
N=500000
\end{lstlisting}

\paragraph{Step 1 — Create pseudo-data.}
Applies systematic per-cell distortion with Gaussian noise:
\begin{lstlisting}
root -l -b -q \
  "create_pseudodata.C(\"$INPUT\",\"$OUTDIR/pseudo.root\",$N,$LEVEL)"
\end{lstlisting}

\paragraph{Step 2 — Derive corrections.}
Compares MC and pseudo-data cell fractions:
\begin{lstlisting}
root -l -b -q \
  "derive_corrections.C(\"$INPUT\",\"$OUTDIR/pseudo.root\",\
  \"$OUTDIR/corrections.root\",$N)"
\end{lstlisting}

\paragraph{Step 3 — Apply Method~1 (flat shift).}
\begin{lstlisting}
root -l -b -q \
  "apply_corrections.C(\"$INPUT\",\"$OUTDIR/corrections.root\",\
  \"$OUTDIR/reweighted_m1.root\",$N,1)"
\end{lstlisting}

\paragraph{Step 4 — Apply Method~2 (TProfile).}
\begin{lstlisting}
root -l -b -q \
  "apply_corrections.C(\"$INPUT\",\"$OUTDIR/corrections.root\",\
  \"$OUTDIR/reweighted_m2.root\",$N,2)"
\end{lstlisting}

\paragraph{Step 5 — Validate (histograms + $\chi^2$).}
\begin{lstlisting}
root -l -b -q \
  "validate_closure.C(\"$INPUT\",\"$OUTDIR\",\
  \"$OUTDIR/closure_histos.root\",$N)"
\end{lstlisting}

\paragraph{Step 6 — Plot.}
\begin{lstlisting}
root -l -b -q \
  "plot_closure.C(\"$OUTDIR/closure_histos.root\",\"$OUTDIR/\")"
\end{lstlisting}

% ---------- 6.3 Closure test: distortion model -----------------------------
\subsubsection{Distortion model}

The pseudo-data generation (\texttt{create\_pseudodata.C}) applies the
following transformation to each cell energy $e_k$:
\begin{equation}
  e'_k = e_k \cdot \bigl(1 + \delta \cdot b_k + \mathcal{G}(0,\sigma)\bigr),
  \label{eq:distortion}
\end{equation}
where $\delta$ is the distortion level, $b_k$ is a fixed bias pattern
based on the Chebyshev distance from the central cell (\cref{tab:bias}),
and $\mathcal{G}(0,\sigma)$ is per-event per-cell Gaussian noise with
$\sigma = \delta/2$ by default. After scaling, total cluster energy is
conserved:
\begin{equation}
  e'_k \leftarrow e'_k \cdot \frac{E_\text{tot}}{\sum_j e'_j}.
\end{equation}

\begin{table}[htbp]
\centering
\caption{Bias pattern $b_k$ as a function of Chebyshev distance from
  the central cell (cell~38).}
\label{tab:bias}
\begin{tabular}{lcl}
\toprule
\textbf{Distance} & \textbf{$b_k$} & \textbf{Region} \\
\midrule
0 & $-1.0$ & Central cell \\
1 & $-0.4$ & First ring (8 cells) \\
2 & $+0.4$ & Second ring (16 cells) \\
$\geq 3$ & $+1.0$ & Outer cells \\
\bottomrule
\end{tabular}
\end{table}

This pattern suppresses the cluster centre and enhances the tails,
producing broader showers that mimic the standard ATLAS data--MC
difference (MC has narrower showers than data).

% ---------- 6.4 Closure test: correction methods ----------------------------
\subsubsection{Correction methods}
\label{sec:running:methods}

Cell energy fractions are defined as $f_k = e_k / E_\text{tot}$ where
$E_\text{tot} = \sum_j e_j$.

\paragraph{Method~1 — Flat shift.}
A single additive correction per cell per $|\eta|$ bin:
\begin{equation}
  \Delta_k = \langle f_k^\text{pseudo}\rangle
            - \langle f_k^\text{MC}\rangle.
\end{equation}
Application: $\;f_k^\text{corr} = f_k^\text{MC} + \Delta_k$.

\paragraph{Method~2 — Energy-dependent TProfile correction.}
A \texttt{TProfile} records the mean correction as a function of the MC
cell fraction:
\begin{equation}
  \alpha_k(f) = \bigl\langle f_k^\text{pseudo} - f_k^\text{MC}
  \bigr\rangle\Big|_{f_k^\text{MC} = f}.
\end{equation}
Application: $\;f_k^\text{corr} = f_k^\text{MC} +
\alpha_k\!\bigl(f_k^\text{MC}\bigr)$, where $\alpha_k$ is obtained
via \texttt{TProfile::Interpolate()}.

In both methods, cluster energy is conserved by rescaling:
\begin{equation}
  e_k^\text{corr} = \frac{f_k^\text{corr}}{\sum_j f_j^\text{corr}}
  \cdot E_\text{tot}.
\end{equation}

% ---------- 6.5 Closure test: quality cuts ----------------------------------
\subsubsection{Quality cuts}
\label{sec:running:quality}

Two quality requirements are applied consistently across all pipeline
stages:

\begin{enumerate}[leftmargin=2em]
  \item \textbf{Negative energy clamping}
    (\texttt{clampCellEnergies()}): Cell energies below zero (from
    electronic noise subtraction) are set to zero before computing
    fractions.

  \item \textbf{Central cell dominance}
    (\texttt{isHealthyCluster()}): The central cell (index~38) must
    have the highest energy fraction. This rejects mis-centred clusters.
\end{enumerate}

% ---------- 6.6 Configurable parameters in config.h ------------------------
\subsubsection{Configuration parameters (\texttt{config.h})}
\label{sec:running:config}

Key constants defined in \texttt{config.h}:

\begin{table}[htbp]
\centering
\caption{Configurable parameters in \texttt{config.h}.}
\label{tab:config}
\small
\begin{tabular}{llp{7cm}}
\toprule
\textbf{Constant} & \textbf{Value} & \textbf{Description} \\
\midrule
\texttt{kEtaSize}, \texttt{kPhiSize}
  & 7, 11 & Cluster dimensions ($\eta\times\phi$) \\
\texttt{kCentralCell} & 38 & Flat index of the central cell \\
\texttt{kDeltaEta} & 0.025 & Layer-2 cell $\eta$ granularity \\
\texttt{kNEtaBins} & 14 & Number of $|\eta|$ bins \\
\texttt{kEtaLimits[15]}
  & $\{0, 0.2, \ldots, 2.4\}$
  & $|\eta|$ bin edges \\
\texttt{kNFracBins} & 12 & TProfile energy-fraction bins \\
\texttt{kFracBins[13]}
  & $\{-1, -0.5, \ldots, 0.5, 1\}$
  & Energy-fraction bin edges \\
\texttt{kTreeName} & \texttt{"tree"} & TTree name in ntuple files \\
\texttt{kCellBranch}
  & \texttt{"photon.7x11ClusterLr2E"}
  & Cell energy branch name \\
\texttt{kEtaBranch}
  & \texttt{"photon\_cluster.eta2"}
  & Cluster $\eta$ branch name \\
\bottomrule
\end{tabular}
\end{table}

% ---------- 6.7 Grid submission --------------------------------------------
\subsection{Grid submission (ntuple production)}
\label{sec:running:grid}

Ntuples are produced from AOD files using the NTupleMaker framework
submitted via pathena. A valid grid proxy is required.

\begin{lstlisting}
# Set up grid tools
setupATLAS
asetup Athena,25.0.40
source NTupleMaker_workspace/build/x86_64-el9-gcc14-opt/setup.sh
lsetup panda rucio
voms-proxy-init -voms atlas

# Submit MC23e Z->llgamma jobs
cd NTupleMaker_workspace/run/
bash ShowerShapesAnalysis/grid/submit_mc23e_Zllg.sh
\end{lstlisting}

The grid scripts automatically set the correct output dataset names,
number of files per job, and analysis type (\texttt{eegamma} or
\texttt{mumugamma}).

% ============================================================================
\section{Output}
\label{sec:output}
% ============================================================================

All output is written under the \texttt{output/} directory and is
gitignored.

\subsection{Validation output}

\begin{table}[htbp]
\centering
\caption{Output files from the validation workflow.}
\label{tab:output:validation}
\small
\begin{tabular}{lp{8cm}}
\toprule
\textbf{File} & \textbf{Contents} \\
\midrule
\texttt{output/reta\_Zeeg.root}
  & TH1F histograms: \texttt{h\_reta\_fudged},
    \texttt{h\_reta\_unfudged}, \texttt{h\_reta\_computed},
    difference histograms, and a 2D correlation histogram. \\
\texttt{output/rphi\_Zeeg.root}
  & Same structure for $R_\phi$. \\
\texttt{output/weta2\_Zeeg.root}
  & Same structure for $w_{\eta\,2}$. \\
\texttt{output/plots/reta\_Zeeg.pdf}
  & Single-page PDF: overlay + ratio plot for $R_\eta$. \\
\texttt{output/plots/rphi\_Zeeg.pdf}
  & Same for $R_\phi$. \\
\texttt{output/plots/weta2\_Zeeg.pdf}
  & Same for $w_{\eta\,2}$. \\
\bottomrule
\end{tabular}
\end{table}

Identical files are produced for each sample (\texttt{Zeeg},
\texttt{Zmumug}).

\subsection{Closure test output}

For each distortion level, a directory is created:

\begin{lstlisting}
output/closure_test_<N>pct/     # N = 1, 3, or 5
\end{lstlisting}

\begin{table}[htbp]
\centering
\caption{Output files per distortion level.}
\label{tab:output:closure}
\small
\begin{tabular}{lp{8cm}}
\toprule
\textbf{File} & \textbf{Contents} \\
\midrule
\texttt{pseudo.root}
  & TTree with distorted cell energies (pseudo-data). \\
\texttt{corrections.root}
  & Method~1: TH2D per $|\eta|$ bin
    (\texttt{flatShift\_eta\_X.XX\_X.XX}).
    Method~2: 77~TProfiles per $|\eta|$ bin
    (\texttt{CellCorrection\_Cell\_K\_Inc\_Eta\_X.XX\_X.XX}). \\
\texttt{reweighted\_m1.root}
  & TTree with Method~1 corrected cell energies. \\
\texttt{reweighted\_m2.root}
  & TTree with Method~2 corrected cell energies. \\
\texttt{closure\_histos.root}
  & TH1D histograms per $|\eta|$ bin per variable per source
    (MC, pseudo, M1, M2). Naming:
    \texttt{h\_<var>\_<source>\_eta<N>}. \\
\texttt{closure\_reta.pdf}
  & Multi-page PDF (13~pages): $R_\eta$ closure plots,
    one page per populated $|\eta|$ bin. \\
\texttt{closure\_rphi.pdf}
  & Same for $R_\phi$. \\
\texttt{closure\_weta2.pdf}
  & Same for $w_{\eta\,2}$. \\
\bottomrule
\end{tabular}
\end{table}

\paragraph{Histogram naming convention.}
Within \texttt{closure\_histos.root}, histograms follow the pattern:
\begin{lstlisting}
h_<variable>_<source>_eta<bin>
\end{lstlisting}
where \texttt{<variable>} $\in$ \{\texttt{reta}, \texttt{rphi},
\texttt{weta2}\}, \texttt{<source>} $\in$ \{\texttt{mc},
\texttt{pseudo}, \texttt{m1}, \texttt{m2}\}, and
\texttt{<bin>} $\in [0, 13]$.

\paragraph{Plot format.}
Each closure plot page contains:
\begin{itemize}[leftmargin=2em]
  \item \textbf{Upper panel}: Normalised overlay of pseudo-data
    (black markers), MC (red filled), Method~1 (green dashed),
    Method~2 (blue solid).
  \item \textbf{Lower panel}: Ratio to pseudo-data (range 0.8--1.2).
  \item \textbf{Text labels}: $|\eta|$ bin range, distortion level.
\end{itemize}

% ============================================================================
\section{Reproducibility}
\label{sec:reproducibility}
% ============================================================================

\subsection{Full reproduction from scratch}

To reproduce the complete set of results, execute the following
commands in order:

\begin{lstlisting}
# 1. Set up environment
setupATLAS
lsetup "root 6.32.08-x86_64-el9-gcc13-opt"

# 2. Navigate to the repository
cd /path/to/ShowerShapesAnalysis

# 3. Ensure input ntuple is present
ls -lh ntuples/mc23e/mc23e_700770_Zeeg.root
# Expected: ~13 GB, ~9.6M events

# 4. Run validation (all 3 variables, both channels)
cd scripts
mkdir -p ../output/plots

python compute_reta.py   -i ../ntuples/mc23e/mc23e_700770_Zeeg.root \
                          -o ../output/reta_Zeeg.root
python compute_rphi.py   -i ../ntuples/mc23e/mc23e_700770_Zeeg.root \
                          -o ../output/rphi_Zeeg.root
python compute_weta_2.py -i ../ntuples/mc23e/mc23e_700770_Zeeg.root \
                          -o ../output/weta2_Zeeg.root

python compute_reta.py   -i ../ntuples/mc23e/mc23e_700771_Zmumug.root \
                          -o ../output/reta_Zmumug.root
python compute_rphi.py   -i ../ntuples/mc23e/mc23e_700771_Zmumug.root \
                          -o ../output/rphi_Zmumug.root
python compute_weta_2.py -i ../ntuples/mc23e/mc23e_700771_Zmumug.root \
                          -o ../output/weta2_Zmumug.root

root -l -b -q 'plot_reta.C("../output/reta_Zeeg.root",\
  "../output/plots/reta_Zeeg")'
root -l -b -q 'plot_rphi.C("../output/rphi_Zeeg.root",\
  "../output/plots/rphi_Zeeg")'
root -l -b -q 'plot_weta_2.C("../output/weta2_Zeeg.root",\
  "../output/plots/weta2_Zeeg")'

root -l -b -q 'plot_reta.C("../output/reta_Zmumug.root",\
  "../output/plots/reta_Zmumug")'
root -l -b -q 'plot_rphi.C("../output/rphi_Zmumug.root",\
  "../output/plots/rphi_Zmumug")'
root -l -b -q 'plot_weta_2.C("../output/weta2_Zmumug.root",\
  "../output/plots/weta2_Zmumug")'

# 5. Run closure test suite (all 3 distortion levels)
cd closure_test
bash run_closure_suite.sh -1    # -1 = all events (~9.6M)
\end{lstlisting}

\subsection{Expected runtime}

\begin{table}[htbp]
\centering
\caption{Approximate runtime on lxplus (single core).}
\label{tab:runtime}
\begin{tabular}{lr}
\toprule
\textbf{Task} & \textbf{Time} \\
\midrule
Validation (1 variable, 1 sample, 9.6M events) & $\sim$15 min \\
Validation (all 3 variables, 2 samples) & $\sim$90 min \\
Closure test (1~level, 9.6M events, all 6 steps) & $\sim$60 min \\
Closure test suite (3 levels, 9.6M events) & $\sim$3 hours \\
Closure test (500k events, quick test) & $\sim$5 min \\
\bottomrule
\end{tabular}
\end{table}

\subsection{Random seeds}

The pseudo-data generation uses \texttt{TRandom3} with seed~0 (which
initialises from the system clock UUID). This means that pseudo-data
files are \emph{not} bitwise-reproducible across runs by default. To
obtain deterministic results, modify \texttt{create\_pseudodata.C} to
use a fixed seed:

\begin{lstlisting}[style=cppstyle]
TRandom3 rng(42);  // Replace rng(0) with fixed seed
\end{lstlisting}

The $\chi^2$ closure metrics are stable across random realisations
because they are averaged over $\sim$7.8M events (after quality cuts).

\subsection{Version control}

The repository tracks only source code and configuration. All input
data (ntuples), output files (histograms, plots, ROOT files), and
compiled artifacts are excluded via \texttt{.gitignore}.

Key gitignore rules:
\begin{lstlisting}
ntuples/           # Input ntuples (large)
output/            # All computed outputs
*.root             # ROOT files in top directory
__pycache__/       # Python bytecode
*.log              # Log files
\end{lstlisting}

\paragraph{Recommended commit workflow.}
\begin{lstlisting}
# Check what changed
git status
git diff --stat

# Stage and commit with conventional prefix
git add -A
git commit -m "feat: add new shower shape variable"
git push origin main
\end{lstlisting}

Use conventional commit prefixes: \texttt{feat:} (new feature),
\texttt{fix:} (bug fix), \texttt{docs:} (documentation),
\texttt{chore:} (maintenance), \texttt{refactor:} (restructuring).

\end{document}
